%% ============================================================
%%  Section 6 — Conclusions
%%  Article: BRKGA with Rolling Horizon for the MDRP
%%  Citation style: LNCS (numeric \cite{})
%% ============================================================

\section{Conclusions}\label{sec:conclusions}

This paper compared four rolling-horizon algorithms for the
Meal Delivery Routing Problem on 24 real-world instances from
a Latin American food-delivery platform.
The key findings are as follows.

\textbf{Coverage quality.}
BRKGA-Coverage and BRKGA-Travel achieve nearly identical
average coverage rates (59.2\%), outperforming the Business Model
Lookahead (51.4\%) and the plain Business Model (44.1\%) on
Small and Medium instances.
On Large instances ($|O|>20{,}000$) the ranking reverses:
the Business Model Lookahead reaches 51.1\% versus 47.2\% for
BRKGA-Coverage, demonstrating that deterministic greedy methods
with demand-forecast deferral scale more gracefully when the
search budget per window is insufficient for effective
evolutionary exploration.

\textbf{Travel efficiency.}
Average travel time per delivered order is virtually identical
across both BRKGA variants (27.9 min/order), confirming that
the nearest-neighbour decoder dominates route quality regardless
of the fitness function used.

\textbf{Computational cost.}
Both BRKGA variants are faster in aggregate run time than the
Business Model Lookahead (138 s and 105 s vs 589 s on average),
because the bounded per-window generation limit keeps each
window's footprint small.
BRKGA-Travel is the fastest overall, 24\% quicker than
BRKGA-Coverage while delivering equivalent coverage.

\textbf{Practical recommendation.}
For small-to-medium instances, BRKGA-Travel offers the best
balance of coverage quality and speed.
For large instances, the Business Model Lookahead is the
preferred choice due to its superior coverage and predictable
scaling behaviour.
A hybrid deployment strategy — switching between the two
based on daily order volume — can capture the best of both worlds.

\textbf{Future work.}
Promising directions include adaptive generation limits that
expand the BRKGA's budget on large instances, incorporation of
courier fairness objectives into the rolling-horizon framework,
and online learning of the demand-forecast parameters from
historical platform data.
